\documentclass[10pt,a4paper,twoside,bg=print,twocolumn,openany,nodeprecatedcode]{dndbook}
\usepackage[utf8]{inputenc}
\usepackage{tabulary}
\usepackage[normalem]{ulem}
\usepackage{color,soul}
\usepackage{float}
\usepackage{caption}
\usepackage{wrapfig}
\usepackage{tocloft}
\usepackage[toc]{multitoc}
\usepackage[hidelinks]{hyperref}
\raggedbottom
\begin{document}
\mainmatter

\DndSpellHeader%
{Cone of Cold}
{5th level Evocation}
{1 action}
{Self (60 feet cone)}
{V, S, a small crystal or glass cone}
{Instantaneous}
You unleash a blast of cold air. Each creature in a 60-foot Cone originating from you makes a Constitution saving throw, taking 8d8 Cold damage on a failed save or half as much damage on a successful one. A creature killed by this spell becomes a frozen statue until it thaws.

\addcontentsline{toc}{section}{Fireball}


\DndSpellHeader%
{Fireball}
{3rd level Evocation}
{1 action}
{150 feet}
{V, S, a ball of bat guano and sulfur}
{Instantaneous}
A bright streak flashes from you to a point you choose within range and then blossoms with a low roar into a fiery explosion. Each creature in a 20-foot-radius Sphere centered on that point makes a Dexterity saving throw, taking 8d6 Fire damage on a failed save or half as much damage on a successful one.

Flammable objects in the area that aren't being worn or carried start burning.

\addcontentsline{toc}{section}{Fog Cloud}


\DndSpellHeader%
{Fog Cloud}
{1st level Conjuration}
{1 action}
{120 feet}
{V, S}
{Concentration, up to 1 hour}
You create a 20-foot-radius Sphere of fog centered on a point within range. The Sphere is Heavily Obscured. It lasts for the duration or until a strong wind (such as one created by \textit{Gust of Wind}) disperses it.

\addcontentsline{toc}{section}{Globe of Invulnerability}


\DndSpellHeader%
{Globe of Invulnerability}
{6th level Abjuration}
{1 action}
{Self (10 feet emanation)}
{V, S, a glass bead}
{Concentration, up to 1 minute}
An immobile, shimmering barrier appears in a 10-foot Emanation around you and remains for the duration.

Any spell of level 5 or lower cast from outside the barrier can't affect anything within it. Such a spell can target creatures and objects within the barrier, but the spell has no effect on them. Similarly, the area within the barrier is excluded from areas of effect created by such spells.

\addcontentsline{toc}{section}{Hold Monster}


\DndSpellHeader%
{Hold Monster}
{5th level Enchantment}
{1 action}
{90 feet}
{V, S, a straight piece of iron}
{Concentration, up to 1 minute}
Choose a creature that you can see within range. The target must succeed on a Wisdom saving throw or have the Paralyzed condition for the duration. At the end of each of its turns, the target repeats the save, ending the spell on itself on a success.

\addcontentsline{toc}{section}{Ice Storm}


\DndSpellHeader%
{Ice Storm}
{4th level Evocation}
{1 action}
{300 feet}
{V, S, a mitten}
{Instantaneous}
Hail falls in a 20-foot-radius, 40-foot-high Cylinder centered on a point within range. Each creature in the Cylinder makes a Dexterity saving throw. A creature takes 2d10 Bludgeoning damage and 4d6 Cold damage on a failed save or half as much damage on a successful one.

Hailstones turn ground in the Cylinder into Difficult Terrain until the end of your next turn.

\addcontentsline{toc}{section}{Levitate}


\DndSpellHeader%
{Levitate}
{2nd level Transmutation}
{1 action}
{60 feet}
{V, S, a metal spring}
{Concentration, up to 10 minutes}
One creature or loose object of your choice that you can see within range rises vertically up to 20 feet and remains suspended there for the duration. The spell can levitate an object that weighs up to 500 pounds. An unwilling creature that succeeds on a Constitution saving throw is unaffected.

The target can move only by pushing or pulling against a fixed object or surface within reach (such as a wall or a ceiling), which allows it to move as if it were climbing. You can change the target's altitude by up to 20 feet in either direction on your turn. If you are the target, you can move up or down as part of your move. Otherwise, you can take a Magic action to move the target, which must remain within the spell's range.

When the spell ends, the target floats gently to the ground if it is still aloft.

\addcontentsline{toc}{section}{Lightning Bolt}


\DndSpellHeader%
{Lightning Bolt}
{3rd level Evocation}
{1 action}
{Self (100 feet line)}
{V, S, a bit of fur and a crystal rod}
{Instantaneous}
A stroke of lightning forming a 100-foot-long, 5-foot-wide Line blasts out from you in a direction you choose. Each creature in the Line makes a Dexterity saving throw, taking 8d6 Lightning damage on a failed save or half as much damage on a successful one.

\addcontentsline{toc}{section}{Magic Missile}


\DndSpellHeader%
{Magic Missile}
{1st level Evocation}
{1 action}
{120 feet}
{V, S}
{Instantaneous}
You create three glowing darts of magical force. Each dart strikes a creature of your choice that you can see within range. A dart deals 1d4 + 1 Force damage to its target. The darts all strike simultaneously, and you can direct them to hit one creature or several.

\addcontentsline{toc}{section}{Ray of Enfeeblement}


\DndSpellHeader%
{Ray of Enfeeblement}
{2nd level Necromancy}
{1 action}
{60 feet}
{V, S}
{Concentration, up to 1 minute}
A beam of enervating energy shoots from you toward a creature within range. The target must make a Constitution saving throw. On a successful save, the target has Disadvantage on the next attack roll it makes until the start of your next turn.

On a failed save, the target has Disadvantage on Strength-based D20 Test for the duration. During that time, it also subtracts 1d8 from all its damage rolls. The target repeats the save at the end of each of its turns, ending the spell on a success.

\addcontentsline{toc}{section}{Wall of Force}


\DndSpellHeader%
{Wall of Force}
{5th level Evocation}
{1 action}
{120 feet}
{V, S, a shard of glass}
{Concentration, up to 10 minutes}
An Invisible wall of force springs into existence at a point you choose within range. The wall appears in any orientation you choose, as a horizontal or vertical barrier or at an angle. It can be free floating or resting on a solid surface. You can form it into a hemispherical dome or a globe with a radius of up to 10 feet, or you can shape a flat surface made up of ten 10-foot-by-10-foot panels. Each panel must be contiguous with another panel. In any form, the wall is 1/4 inch thick and lasts for the duration. If the wall cuts through a creature's space when it appears, the creature is pushed to one side of the wall (you choose which side).

Nothing can physically pass through the wall. It is immune to all damage and can't be dispelled by \textit{Dispel Magic}. A \textit{Disintegrate} spell destroys the wall instantly, however. The wall also extends into the Ethereal Plane and blocks ethereal travel through the wall.

\addcontentsline{toc}{section}{Wall of Ice}


\DndSpellHeader%
{Wall of Ice}
{6th level Evocation}
{1 action}
{120 feet}
{V, S, a piece of quartz}
{Concentration, up to 10 minutes}
You create a wall of ice on a solid surface within range. You can form it into a hemispherical dome or a globe with a radius of up to 10 feet, or you can shape a flat surface made up of ten 10-foot-square panels. Each panel must be contiguous with another panel. In any form, the wall is 1 foot thick and lasts for the duration.

If the wall cuts through a creature's space when it appears, the creature is pushed to one side of the wall (you choose which side) and makes a Dexterity saving throw, taking 10d6 Cold damage on a failed save or half as much damage on a successful one.

The wall is an object that can be damaged and thus breached. It has AC 12 and 30 Hit Points per 10-foot section, and it has Immunity to Cold, Poison, and Psychic damage and Vulnerability to Fire damage. Reducing a 10-foot section of wall to 0 Hit Points destroys it and leaves behind a sheet of frigid air in the space the wall occupied.

A creature moving through the sheet of frigid air for the first time on a turn makes a Constitution saving throw, taking 5d6 Cold damage on a failed save or half as much damage on a successful one.

\end{document}
